\begin{longtable}{@{}l p{12.5cm}@{}}
  \caption{Skills per-instance answer comparison ($0 < \Delta < 25\%$) for the 15 most divergent items: same question, each model's answer and reasoning.}\label{tab:cmp_skills_pos0} \\
  \toprule
  \endfirsthead
  \toprule
  \endhead
  \multicolumn{2}{@{}l}{\textbf{58191} \quad $\Delta = +20\%$} \\
  \textbf{Question} & You are a tutor conducting a progress check and feedback session for a kindergarten student in an arts unit. The student has completed a portfolio containing: three mixed-media collages exploring 'My Family', a short oral description of each collage (recorded), and a photographed timeline showing their process over four lessons. Using the portfolio and the tutor's notes (which indicate the student used color, texture, and simple symbolic shapes but had difficulty explaining choices and sequencing steps), summarize the student's current level of understanding in art concepts and processes. Diagnose specific gaps in their comprehension and skills related to artistic expression, materials knowledge, and metacognitive reflection. Then design a concrete, age-appropriate improvement plan with three sequential activities (one per week) that scaffold the student's abilities toward clearer expressive intent, better use of materials, and improved ability to describe and sequence their creative process. For each activity, include: objective, materials, step-by-step procedure, expected student response, and one measurable way you will check progress at the end of that week. Be sure the activities are suitable for kindergarteners and support both artistic and language development. \\
  \textbf{Reference} & Summary of current level: \newline The student demonstrates emergent artistic understanding appropriate for kindergarten with strengths in experimentation and basic symbolic representation. In the collages they consistently choose colors and textures to represent family members and include simple symbolic shapes (e.g., circles for faces, rectangles for bodies). Their process photos show they engage in sequential steps (gathering materials, arranging, gluing), and they sustain attention across multiple lessons. However, the student has limited ability to articulate the reasons behind choices (color meaning, material selection) and difficulty verbally sequencing steps beyond immediate actions. Their fine-motor control is developing: collage cutting and gluing are adequate but would benefit from more deliberate layering and control of materials. \newline Diagnosed gaps: \newline 1) Expressive intent: The student uses symbols but cannot consistently explain why a color or texture represents a family member (e.g., "red because mommy is nice" lacks specificity or is missing). This suggests nascent symbolic reasoning and limited vocabulary for describing feelings and attributes. \newline 2) Materials knowledge and technique: The student experiments but shows inconsistent use of materials (over-application of glue, uneven layering, limited exploration of combining textures deliberately). They need guided practice with safe tool control and purposeful material choices. \newline 3) Metacognitive reflection and sequencing: The student struggles to narrate the sequence of their process and reflect on decisions. They can perform steps but cannot reliably tell the order or justify modifications, limiting ability to plan and improve. \newline \textit{[\ldots\ truncated \ldots]} \newline Overall expected growth and next steps: \newline Across the three weeks, the student should show clearer links between color/texture and meaning, better control and purposeful use of materials, and improved ability to plan and narrate a simple three-step process. After Week 3, continue integrating short reflection prompts into each art lesson (e.g., a one-sentence caption) and gradually increase the complexity of planning (4-step storyboards) as language and sequencing skills strengthen. \\
  \midrule
  \textbf{Kimi-K2.6} & score 7/10 \\
  & The user wants a detailed educational response structured as a tutor's progress check and feedback session for a kindergarten student in an arts unit. I need to: \newline 1. Summarize the student's current level of understanding in art concepts and processes based on the portfolio (3 mixed-media collages on 'My Family', oral descriptions recorded, photographed timeline over 4 lessons) and tutor notes (used color, texture, simple symbolic shapes; difficulty explaining choices and sequencing steps). \newline 2. Diagnose specific gaps in comprehension and skills related to: \newline - Artistic expression \newline - Materials knowledge \newline - Metacognitive reflection \newline \textit{[\ldots\ truncated \ldots]} \newline - Metacognitive reflection: Difficulty with causal language ("because") and temporal sequencing ("first/then/last"); struggles to externalize thinking. \newline The plan looks good. I'll format it cleanly with headers and bullet points for readability. \\
  \textbf{Innospark-1T} & score 9/10 \\
  & Let me think through this carefully. \newline The question/task: You are a tutor conducting a progress check and feedback session for a kindergarten student in an arts unit. The student has completed a portfolio containing: three mixed-media collages exploring 'My Family', a short oral description of each collage (recorded), and a photographed timeline showing their process over four lessons. Using the portfolio and the t \newline My reasoning: Progress Check and Feedback Session: Kindergarten Arts Unit \newline Student Portfolio Analysis \newline Summary of Current Understanding \newline Based on the portfolio evidence and tutor notes, this kindergarten student demonstrates: \newline \textit{[\ldots\ truncated \ldots]} \newline | Week 3 | Sequenced reflection | Creates coherent process narrative | \newline This scaffolded approach builds from basic expressive language (Week 1), through intentional material use (Week 2), to integrated metacognitive reflection (Week 3), supporting both artistic and language development appropriate for kindergarten learners. \\
  \specialrule{1.2pt}{8pt}{8pt}
  \multicolumn{2}{@{}l}{\textbf{19775} \quad $\Delta = +20\%$} \\
  \textbf{Question} & Using Gagn's ABCD (Audience, Behavior, Condition, Degree) framework and the results of a task analysis for an advanced undergraduate physics course module on electromagnetic wave propagation in anisotropic media, design an assessment item that validly measures whether students have internalized the value of modeling assumptions when applying Maxwell's equations to real-world materials. Your response should (1) state a precise learning objective in ABCD format including the specific learning content, (2) describe the assessment task you would administer (format, conditions, and degree of mastery required), and (3) justify how the task aligns with the task analysis outcomes and promotes characterizing-level internalization of valuing appropriate modeling assumptions (i.e., students consistently prioritize, critique, and choose modeling assumptions in novel contexts). Provide concrete examples of acceptable student responses that demonstrate this internalization. \\
  \textbf{Reference} & 1) Learning objective (ABCD): Given advanced undergraduate students in the Electromagnetism course (Audience), students will evaluate and select appropriate modeling assumptions and apply Maxwell's equations to derive boundary conditions and wave solutions for electromagnetic propagation in a specified anisotropic dielectric slab under realistic measurement constraints (Behavior), using provided material response data, uncertain parameter ranges, and experimental setup limitations (Condition), with justification of chosen assumptions and predicted error bounds such that predicted field amplitudes and phase velocities lie within 10\% of values obtained by a provided high-fidelity numerical model (Degree). Learning content: Maxwell's equations in matter, constitutive relations for anisotropic permittivity tensors, boundary-value problems for stratified media, uncertainty propagation, and critical assessment of modeling assumptions (e.g., linearity, homogeneity, frequency independence, neglect of loss, effective medium approximations). \newline 2) Assessment task (format, conditions, degree): Format: A multi-part open-response task to be completed in 90 minutes, including written derivations, selection and justification of modeling assumptions, and a short computational component (pseudo-code or algorithmic description sufficient; no full coding required). Conditions: Students receive a problem brief describing a laboratory scenario---an anisotropic dielectric slab (uniaxial crystal) of unknown thickness sandwiched between two isotropic dielectrics, measured reflectance and transmittance spectra at several angles, and manufacturer-provided ranges for permittivity tensor components and slab thickness (with 20\% uncertainty). They are told that measurement noise and finite beam divergence are present. Degree of mastery required: Students must (a) explicitly list and prioritize the modeling assumptions they adopt for solving the boundary-value problem, (b) derive the boundary conditions and analytic or semi-analytic expressions for transmitted and reflected fields under those assumptions, (c) use the provided uncertainty ranges to estimate propagated uncertainty in key observables (e.g., transmitted intensity and phase) and show that their predictions match the provided high-fidelity numerical model within 10\%, and (d) critically justify, with reference to the task analysis, why alternative assumptions were rejected and under what conditions those alternatives would become preferable. \newline 3) Justification of alignment with task analysis and promotion of characterizing-level internalization: The task analysis identified component skills: deriving Maxwell-based boundary-value formulations, manipulating anisotropic constitutive tensors, performing uncertainty propagation, interpreting experimental constraints, and exercising metacognitive judgment about model applicability. The assessment integrates those skills into a complex, authentic decision-making task requiring synthesis and value-based prioritization (e.g., choosing which approximations to make when faced with limited or noisy data). By forcing students to both make and defend modeling choices while quantifying their consequences relative to a high-fidelity reference, the task measures not merely technical skill but the internalized disposition to privilege physically justified, uncertainty-aware assumptions in new contexts --- the hallmark of characterizing. \newline 4) Concrete examples of acceptable student responses demonstrating internalization: Example response A (exemplary): "Assumptions I adopt, ranked: (1) Treat the uniaxial slab as locally homogeneous at the scale of the wavelength because the manufacturer specifies grain sizes \textless{}\textless{} wavelength; (2) Retain anisotropic, frequency-dependent permittivity tensor but linear response; (3) Include small but nonzero loss (imaginary parts) because measured transmittance dips indicate absorption; (4) Model the incident beam as a finite-angle Gaussian rather than a plane wave to account for angular averaging in measurements; (5) Neglect magnetization (0) given non-magnetic materials. Rationale: Homogeneity simplifies boundary conditions and is supported by task analysis noting spatial scales; retaining anisotropy preserves essential physics; including loss is required to match observed attenuation. \newline Using these assumptions, I write Maxwell's equations with D = E and impose tangential E and H continuity at interfaces. For TE and TM polarizations I derive dispersion relations for ordinary and extraordinary waves in the uniaxial slab. Applying Fresnel-like matrix methods for anisotropic stratified media yields expressions for complex reflection and transmission coefficients R(,) and T(,). To account for beam divergence, I integrate the angular-dependent T(,) weighted by a Gaussian angular spectrum with corresponding to the documented beam profile. \newline Uncertainty propagation: I perform first-order (linear) propagation of parameter uncertainties: T (T/p\_i) p\_i where p\_i are \_xx, \_zz, and slab thickness d. Using manufacturer 20\% bounds and measured noise, propagated uncertainty yields predicted transmitted intensity I\_pred 8\% at the principal measurement frequencies. Comparison with the provided high-fidelity model shows agreement within 6\% across the central frequencies. \newline \textit{[\ldots\ truncated \ldots]} \newline Example response B (satisfactory): "I assume local homogeneity, linear anisotropic permittivity with small loss, and plane-wave incidence but correct results using an angular correction factor. I derive boundary conditions and use transfer-matrix formalism for anisotropic layers to compute T(,). Linear uncertainty propagation with the given 20\% parameter ranges gives predicted intensities within 9\%, matching the reference model. I note that isotropic approximations fail to reproduce polarization splitting, and including loss is necessary to match attenuation. Thus, my assumptions are justified for this experimental context; if grain size were comparable to wavelength, I would adopt a different approach (homogenization invalid)." \newline Both examples show the student consistently justifying and prioritizing modeling assumptions, quantifying their impact, and describing when choices would change---demonstrating characterizing-level internalization of valuing appropriate modeling assumptions in application of Maxwell's equations to real materials. \\
  \midrule
  \textbf{Kimi-K2.6} & score 6/10 \\
  & [5-6] The response provides a strong ABCD objective, detailed task format, and solid justification, but it is incomplete. The 'Concrete Examples of Acceptable Responses' section cuts off abruptly, omitting the actual examples. This makes the scoring criteria and correct answers insufficient, requiring revision before practical application. \\
  \textbf{Innospark-1T} & score 8/10 \\
  & Let me think through this carefully. \newline The question/task: Using Gagn's ABCD (Audience, Behavior, Condition, Degree) framework and the results of a task analysis for an advanced undergraduate physics course module on electromagnetic wave propagation in anisotropic media, design an assessment item that validly measures whether students have internalized the value of modeling assumptions when applying Max \newline My reasoning: Assessment Design: Modeling Assumptions in Anisotropic Media \newline Task Analysis Summary \newline Before designing the assessment, I conducted a task analysis identifying the hierarchical skills required for electromagnetic wave propagation in anisotropic media: \newline | Task Level | Component Skills | \newline \textit{[\ldots\ truncated \ldots]} \newline Rubric specifies measurable degree of mastery \newline Examples demonstrate characterizing-level valuing (consistent prioritization, critique, transfer) \\
  \specialrule{1.2pt}{8pt}{8pt}
  \multicolumn{2}{@{}l}{\textbf{114398} \quad $\Delta = +20\%$} \\
  \textbf{Question} & You are a research assistant tasked with generating a prioritized list of novel, high-impact research ideas in engineering education for higher education instructors. Based on an automated literature scan that summarizes recent papers (past 5 years) identifying trends such as active learning adoption, AI-assisted tutoring, competency-based assessment, equity and inclusion, and industry-academia partnerships --- and noting gaps like limited longitudinal studies, few cross-cultural comparisons, lack of implementation-ready intervention packages, and sparse cost-effectiveness analyses --- evaluate and synthesize these findings to: (1) propose five concrete research ideas that address the most critical gaps, (2) justify why each idea is high-priority by linking it to evidence from the summarized trends and gaps, (3) outline for each idea a feasible study design (including population, methods, key measures, and an estimated timeframe), and (4) identify the primary potential barriers to implementation and how you would mitigate them. Provide clear, evidence-linked reasoning for prioritization and feasibility for each proposed idea. \\
  \textbf{Reference} & Model answer: \newline Overview and prioritization rationale \newline Given the literature scan showing widespread interest in active learning, AI-assisted tutoring, competency-based assessment, equity and inclusion, and industry partnerships --- alongside clear gaps in longitudinal evidence, cross-cultural generalizability, turnkey interventions, and cost-effectiveness studies --- I prioritize research that (a) tests implementation-ready solutions, (b) produces generalizable longitudinal evidence, and (c) provides decision-relevant metrics (learning, equity, and cost). The five proposed studies below are ordered by priority based on potential impact, feasibility, and ability to fill multiple gaps simultaneously. \newline 1) Multi-institutional longitudinal evaluation of an implementation-ready active-learning package \newline - Why high priority: Active learning is widely recommended but adoption and sustained effectiveness across contexts are uncertain; longitudinal and multi-institutional evidence is scarce. An implementation-ready package addresses the gap in turnkey interventions. \newline - Study design: Cluster quasi-experimental longitudinal design across 8 universities (mix of research-intensive and teaching-focused, two countries if possible). Population: undergraduate engineering cohorts in core courses (e.g., statics, circuits) across 2 academic years plus follow-up one year post-course. Intervention: a standardized instructor toolkit (syllabus-aligned active-learning modules, assessment templates, faculty training modules). Methods: baseline measurements year 0, implementation in year 1 with matched control sections using propensity matching, follow-up in year 2 to assess retention and transfer. Key measures: concept mastery (validated concept inventories), course grades, retention in program, student engagement surveys, instructor fidelity logs. Timeframe: 3 years (1 year prep, 1 year implementation, 1 year follow-up). \newline \textit{[\ldots\ truncated \ldots]} \newline Conclusion \newline These five studies together address the identified gaps by producing longitudinal, cross-cultural, implementation-ready, and cost-informed evidence. Each proposed design balances rigor and feasibility, anticipates major barriers, and includes concrete mitigation strategies to enhance the likelihood of producing actionable, generalizable findings for higher education engineering contexts. \\
  \midrule
  \textbf{Kimi-K2.6} & score 7/10 \\
  & [7-8] The student provides 3 highly relevant, creative, and detailed research ideas with strong study designs, measures, and mitigations. However, the prompt requested 5 ideas, and the third idea is cut off mid-sentence. This incompleteness reduces the topic diversity and leaves the design details for the final ideas missing, placing it in the 7-8 band. \\
  \textbf{Innospark-1T} & score 9/10 \\
  & Let me think through this carefully. \newline The question/task: You are a research assistant tasked with generating a prioritized list of novel, high-impact research ideas in engineering education for higher education instructors. Based on an automated literature scan that summarizes recent papers (past 5 years) identifying trends such as active learning adoption, AI-assisted tutoring, competency-based assessment, equity and inclusion, and industry-academia p \newline My reasoning: Research Ideas in Engineering Education \newline Analysis of Trends and Gaps \newline Based on the literature scan, I'll first synthesize the key findings: \newline | Trend | Current State | Key Limitation | \newline \textit{[\ldots\ truncated \ldots]} \newline 3. It has the highest feasibility with clear, actionable outcomes \newline This should be followed by Idea 1 (Longitudinal Study) as a complementary effort requiring longer-term commitment but yielding foundational evidence for the field. \\
  \specialrule{1.2pt}{8pt}{8pt}
  \multicolumn{2}{@{}l}{\textbf{42370} \quad $\Delta = +20\%$} \\
  \textbf{Question} & A high school student submitted a written reflection about a PE unit on cardiovascular fitness that contains factual errors: they claimed that "all aerobic exercises are low intensity and only improve muscular strength," and they also wrote that "resting heart rate increases with improved fitness." As a teaching assistant responsible for detecting and correcting assignment errors, identify and explain the factual errors in the student's statements and provide corrected explanations for each point. In your answer, clearly state the incorrect claim, why it is incorrect (briefly), and the accurate statement with a simple supporting fact or definition. \\
  \textbf{Reference} & Incorrect claim 1: "All aerobic exercises are low intensity and only improve muscular strength." Why this is incorrect: Aerobic exercise refers to activities that depend primarily on oxygen to produce energy and can vary widely in intensity; they are not limited to low intensity, nor is their primary effect to improve muscular strength. Accurate statement 1: Aerobic (cardiorespiratory) exercises include both low- and moderate- to high-intensity activities (for example, walking is typically low intensity, while running or high-intensity interval training can be high intensity). The primary adaptations from regular aerobic exercise are improvements in cardiovascular and respiratory systems and endurance (e.g., increased stroke volume, improved oxygen delivery, and enhanced aerobic capacity), while muscular strength is more directly improved by resistance or strength training that uses high-load, low-repetition muscle contractions. Supporting fact: Aerobic fitness is commonly measured by VO2 max or sustained heart rate response, not by one-repetition maximum strength. \newline Incorrect claim 2: "Resting heart rate increases with improved fitness." Why this is incorrect: Improved cardiovascular fitness typically leads to a lower resting heart rate, not a higher one, because the heart becomes more efficient at pumping blood. Accurate statement 2: With improved aerobic fitness and cardiovascular conditioning, resting heart rate usually decreases as the heart's stroke volume increases and it can deliver the same cardiac output with fewer beats per minute. Supporting fact: Trained endurance athletes often have resting heart rates lower than average (for example, in the 40s--50s beats per minute), whereas a sedentary person's resting heart rate is commonly higher. \\
  \midrule
  \textbf{Kimi-K2.6} & score 8/10 \\
  & The user wants me to identify and correct two factual errors in a high school student's PE reflection about cardiovascular fitness. \newline Error 1: "all aerobic exercises are low intensity and only improve muscular strength" \newline - Incorrect claim: All aerobic exercises are low intensity and only improve muscular strength. \newline - Why incorrect: Aerobic exercises can range from low to high intensity (e.g., running, swimming, cycling at various intensities). Also, aerobic exercises primarily improve cardiovascular endurance, not muscular strength (which is more the domain of anaerobic/resistance training). \newline - Correct statement: Aerobic exercises vary in intensity and primarily improve cardiovascular endurance. Supporting fact: Aerobic exercise strengthens the heart and improves the body's ability to use oxygen, while muscular strength is typically developed through resistance training. \newline Error 2: "resting heart rate increases with improved fitness" \newline \textit{[\ldots\ truncated \ldots]} \newline **Accurate Statement:** Resting heart rate typically decreases with improved cardiovascular fitness. Supporting fact: A well-conditioned heart is stronger and more efficient, so it pumps more blood per beat; this is why trained individuals and athletes often have resting heart rates in the range of 40--60 beats per minute, compared to the average untrained adult range of 60--100 bpm. \newline This looks good. It's concise, accurate, and follows the requested format. I'll present it cleanly. \\
  \textbf{Innospark-1T} & score 10/10 \\
  & Let me analyze the student's reflection and identify the factual errors: \newline **Error 1: "All aerobic exercises are low intensity and only improve muscular strength"** \newline Breaking this down: \newline - The student claims ALL aerobic exercises are low intensity \newline - The student claims aerobic exercises ONLY improve muscular strength \newline Why this is incorrect: \newline \textit{[\ldots\ truncated \ldots]} \newline - Well-trained athletes may have resting heart rates of 40-60 bpm (vs. average 60-100 bpm) \newline - This is because the heart muscle strengthens and stroke volume increases \\
  \specialrule{1.2pt}{8pt}{8pt}
  \multicolumn{2}{@{}l}{\textbf{30556} \quad $\Delta = +20\%$} \\
  \textbf{Question} & You are creating a text-based learning activity for kindergarteners that uses simple programming concepts (such as sequences, conditionals expressed in plain language, and simple loops represented by repeated actions) to teach social-emotional skills like sharing, recognizing emotions, and turn-taking. Design a single, detailed learning task (a short scenario and accompanying coding-style instructions children can 'run' using plain-language commands) that asks students to solve a social-emotional problem. Then evaluate and justify your design choices: explain why the chosen scenario is appropriate for kindergarten social-emotional development, how the programming-style instructions scaffold cognitive and socio-emotional learning at a high level, how you balanced simplicity and challenge to meet a 'Hard / Evaluating' learning objective, and how you would expect children to demonstrate their understanding through the activity. Finally, propose one specific modification to adapt the task for a child who is advanced in both code-concepts and social skills, and justify that modification. \\
  \textbf{Reference} & Task design (scenario and instructions): \newline Scenario: "At Sunny Class, three friends (Ava, Ben, and Carlos) want to build a tower with colored blocks, but there is only one red block left and all three want it. They need to take turns and decide who will use the red block now and how to share it afterward so everyone feels happy." \newline Plain-language coding-style instructions children can 'run': \newline 1. Look: CHECK who is closest to the red block. (If two or more are equally close, go to step 3.) \newline 2. Ask: SAY, 'May I use the red block for two turns?' and WAIT for a response. \newline 3. Decide: IF someone says 'Yes' then GIVE a thank-you hug or high-five. ELSE if everyone says 'No' then START a fairness rule: each person gets the red block for 1 turn in order. \newline \textit{[\ldots\ truncated \ldots]} \newline - Modification: Introduce a simple variable and optimization problem: provide a set of colored blocks with limited red blocks and an additional constraint---each child has a different priority score (e.g., Ava needs red for her pattern project, Ben only likes blue, Carlos is allergic to blue). Ask the advanced child to DESIGN a short 'program' (3--6 plain-language commands) that assigns turns to maximize overall happiness while being fair, and to justify the rule they created. \newline - Justification: This increases abstraction by adding state (priority scores) and an optimization/evaluation component (maximizing happiness subject to fairness), engaging higher-order synthesis and evaluative reasoning. It also preserves socio-emotional focus because the student must weigh needs and preferences, practice perspective-taking at a deeper level, and justify ethical choices. The task still uses plain-language code so the child applies computational thinking without requiring formal syntax, making it developmentally appropriate while appropriately challenging. \\
  \midrule
  \textbf{Kimi-K2.6} & score 5/10 \\
  & [5-6] The provided scenario and code blocks are creative and well-aligned with kindergarten SEL and coding concepts, but the answer is incomplete. It cuts off mid-sentence, entirely missing the required evaluation, justification of design choices, explanation of the 'Hard/Evaluating' objective, expected demonstrations of understanding, and the advanced modification. This results in an insufficient educational flow and lacks the necessary assessment and difficulty adjustment analysis. \\
  \textbf{Innospark-1T} & score 7/10 \\
  & Let me think through this carefully. \newline The question/task: You are creating a text-based learning activity for kindergarteners that uses simple programming concepts (such as sequences, conditionals expressed in plain language, and simple loops represented by repeated actions) to teach social-emotional skills like sharing, recognizing emotions, and turn-taking. Design a single, detailed learning task (a short scenario and accompanying \newline My reasoning: \textless{}Thinking\textgreater{} \newline Designing an effective learning activity for kindergarteners requires careful consideration of developmental appropriateness, cognitive load, and the integration of multiple learning objectives. \textless{}/Thinking\textgreater{} \newline \textless{}Answer\textgreater{} \newline Learning Task: "The Feelings Robot at Playtime" \newline \textit{[\ldots\ truncated \ldots]} \newline - This modification maintains age-appropriate content while elevating the thinking required \newline \textless{}/Answer\textgreater{} \\
  \specialrule{1.2pt}{8pt}{8pt}
  \multicolumn{2}{@{}l}{\textbf{108627} \quad $\Delta = +20\%$} \\
  \textbf{Question} & You are an evaluator designing a personalized explanation for a kindergarten student's written math answer. The student's response to the prompt 'Draw 5 apples and show 2 of them being eaten. How many apples are left?' is: '5 apples. 2 eaten. 5 - 2 = 2. So 2 left.' Identify and analyze each sentence in the student's answer that contains an error or misunderstanding. For each identified error, explain clearly what the mistake is, what the correct mathematical idea or calculation should be, and provide a short, child-friendly explanation (one or two sentences) that could be used to teach the student the correct concept. \\
  \textbf{Reference} & Sentence-by-sentence analysis: \newline 1) '5 apples.' \newline - Error or misunderstanding: No error. This sentence correctly states the initial number of apples. \newline - Correct idea: The student begins with 5 apples, which matches the prompt. \newline - Child-friendly explanation: 'Yes --- you started with five apples. Good job counting them!' \newline 2) '2 eaten.' \newline \textit{[\ldots\ truncated \ldots]} \newline - Child-friendly explanation: 'Let's count the apples that aren't eaten: one, two, three. There are three apples left.' \newline Summary for the student (brief corrective explanation you can read aloud): 'You started correctly with five apples and noticed that two were eaten. When we take two away from five, we count 1, 2, 3 apples left. So 5 minus 2 equals 3 --- three apples remain.' \\
  \midrule
  \textbf{Kimi-K2.6} & score 8/10 \\
  & [7-8] The student accurately identifies and analyzes the errors in sentences 3 and 4 with clear, child-friendly explanations. However, it omits the sentence-by-sentence analysis of sentences 1 and 2 (noting they are correct), lacking the comprehensive depth of the reference answer. \\
  \textbf{Innospark-1T} & score 10/10 \\
  & Let me work through this problem step by step. \newline First, I need to understand what's being asked: You are an evaluator designing a personalized explanation for a kindergarten student's written math answer. The student's response to the prompt 'Draw 5 apples and show 2 of them being eaten. How many apples are left?' is: '5 apples. 2 eaten. 5 - 2 = 2. So 2 left.' Identify and analyze each sent \newline Key values given: 5, 2, 5, 2, 2 \newline Working through the calculation: Analysis of Student's Math Response \newline Step-by-Step Evaluation \newline Let me carefully examine each sentence in the student's answer: \newline \textit{[\ldots\ truncated \ldots]} \newline "Let's count together! If you have 5 apples and eat 2, put 2 apples in your tummy!" \newline This gives us the final answer: 5 - 2 = 3 \\
  \specialrule{1.2pt}{8pt}{8pt}
  \multicolumn{2}{@{}l}{\textbf{42640} \quad $\Delta = +10\%$} \\
  \textbf{Question} & You are a teaching assistant reviewing a submitted programming assignment from a high-school career education class. The student's Python function is intended to read a list of integers and return a new list containing only the even numbers squared, preserving the original order. The submitted code is below and fails on some test cases. Identify all logical and runtime errors in the code, explain why each error causes incorrect behavior, and provide a corrected version of the function. Finally, describe how you would test the corrected function with at least four distinct test cases (inputs and expected outputs) that cover normal cases and edge cases. \\
  \textbf{Reference} & Model answer: \newline Analysis of submitted code and identified errors: \newline 1) Suppose the submitted code is: \newline def process\_numbers(nums): \newline result = \newline for i in range(len(nums)): \newline \textit{[\ldots\ truncated \ldots]} \newline Reason: test to ensure correct handling of large integer squares. \newline This corrected implementation and the test cases ensure the function meets the specification: it selects even integers, squares them, preserves order, and avoids runtime errors when encountering non-integer items. \\
  \midrule
  \textbf{Kimi-K2.6} & score 7/10 \\
  & [7-8] The student accurately identifies major errors and provides a valid corrected function with appropriate test cases. However, the cause analysis for each error is brief and lacks the evidence-based depth and prioritization required for a higher score, and it misses the robustness issue (non-integer inputs) highlighted in the reference. \\
  \textbf{Innospark-1T} & score 8/10 \\
  & Let me analyze the student's code carefully. Since no code was provided in the prompt, I'll need to work with a typical erroneous implementation that students might submit for this problem. \newline Common errors students make for "even numbers squared": \newline 1. Checking for odd instead of even (num \% 2 != 0) \newline 2. Forgetting to square the number \newline 3. Modifying the original list instead of creating a new one \newline 4. Using integer division // instead of modulo \% \newline \textit{[\ldots\ truncated \ldots]} \newline 7. Type errors (not checking if input is valid) \newline Let me create a comprehensive response that addresses typical errors and provides thorough testing guidance. \\
  \specialrule{1.2pt}{8pt}{8pt}
  \multicolumn{2}{@{}l}{\textbf{81962} \quad $\Delta = +10\%$} \\
  \textbf{Question} & You are an evaluator producing a diagnostic report for an undergraduate chemistry student who submitted a mixed set of short-answer and multiple-choice responses about acid--base equilibria and titration concepts. Based on the learner's incorrect responses, analyze and classify the most likely underlying cause(s) of their misconceptions. In your analysis, (1) identify at least three distinct types of pre-knowledge gaps or affective factors (e.g., misconceptions about definitions, procedural anxiety, low confidence in conceptual reasoning) that could explain their errors; (2) prioritize these causes by explaining which you believe are most influential and why; (3) propose how these prioritized causes should shape the recommended diagnostic summary and next instructional steps for the instructor, emphasizing how to resolve value conflicts (such as the student's preference for rote procedures versus developing conceptual understanding) and promote motivation and responsibility for learning. Provide a coherent, organized response that integrates cognitive and affective aspects and justifies your prioritization. \\
  \textbf{Reference} & Model answer: \newline Analysis and classification of underlying causes \newline 1) Misconception about fundamental definitions and representations (cognitive pre-knowledge gap). The student consistently answered that the equivalence point in titration is always at pH 7 and labeled indicators by color changes without reference to pKa ranges. They also treated conjugate base strength as synonymous with concentration. These errors indicate incomplete or incorrect mental models of acid--base definitions (pKa vs pH, equivalence vs endpoint) and symbol-meaning mapping (what Ka, pKb, pH signify). This is a primary content knowledge deficit. \newline 2) Overreliance on rote procedural strategies and pattern matching (cognitive-affective hybrid). Several incorrect responses used memorized procedures (add strong base until indicator changes) rather than applying Henderson--Hasselbalch or buffer reasoning to predict pH. The student displayed comfort following steps but difficulty transferring principles to slightly novel situations (e.g., titration of weak acid with strong base). This suggests a learned preference for procedure over conceptual analysis---an affective orientation favoring low-effort, low-risk approaches. \newline 3) Low confidence and test anxiety affecting analytical reasoning (affective factor). In short-answer items that required justification, the student gave brief, uncertain statements and avoided exploring alternative answers. This pattern is consistent with anxiety-driven truncation of responses and risk-avoidant answering behavior, reducing their engagement in deeper reasoning even when partial knowledge exists. \newline 4) Misinterpretation of experimental context and data (skill gap). Errors interpreting titration curves---misplacing equivalence and selecting inappropriate indicator---point to limited ability to translate graphical information into chemical conclusions. This may stem from insufficient practice integrating experimental representations with theory. \newline \textit{[\ldots\ truncated \ldots]} \newline Explicitly address the student's apparent preference for rote procedures by framing conceptual work as efficiency-building rather than extra effort. Present short comparisons showing how conceptual reasoning reduces errors and speeds problem solving in novel contexts---demonstrate a one-minute conceptual checklist that simplifies decision making (identify acid/base strength, estimate pH region, choose indicator). Invite the student to commit to a balanced approach: three conceptual-reflection questions per problem before any calculation. This negotiation respects the learner's value for procedural speed while prioritizing conceptual competence. \newline Finally, foster ownership by co-creating measurable, short-term goals (e.g., correctly justify indicator choice in three successive problems) and scheduling brief progress checks. This aligns affective support with cognitive remediation and makes the diagnostic summary actionable for the instructor. \\
  \midrule
  \textbf{Kimi-K2.6} & score 7/10 \\
  & [7-8] The student accurately identifies and classifies three distinct causes with appropriate rationale, and the prioritization is well-justified. However, the response is truncated, leaving the instructional recommendations and resolution of value conflicts incomplete. This aligns with the 7-8 band where parts of the summary and improvement suggestions lack sufficient detail. \\
  \textbf{Innospark-1T} & score 8/10 \\
  & Let me think through this carefully. \newline The question/task: You are an evaluator producing a diagnostic report for an undergraduate chemistry student who submitted a mixed set of short-answer and multiple-choice responses about acid--base equilibria and titration concepts. Based on the learner's incorrect responses, analyze and classify the most likely underlying cause(s) of their misconceptions. In your analysis, (1) identify at least three disti \newline My reasoning: Diagnostic Report: Analysis of Student Misconceptions in Acid-Base Equilibria and Titration \newline Executive Summary \newline Based on the student's incorrect responses regarding acid-base equilibria and titration concepts, I have identified three primary categories of underlying causes. Below is my systematic analysis with prioritized recommendations. Misconception Category 1: Fragmented Procedural Knowledge Without Conceptual Coherence \newline Description: The student appears to apply formulas (Henderson-Hasselbalch equation, pH calculations) mechanically without understanding the chemical meaning behind the variables. For example, they may calculate buffer pH correctly while failing to explain why the buffer resists pH change at the molecular level. Evidence Indicators: \newline \textit{[\ldots\ truncated \ldots]} \newline Conclusion \newline This student's errors stem primarily from a fragmented approach to learning chemistry---prioritizing calculation procedures over conceptual coherence. By addressing the cognitive gaps systematically while simultaneously reducing anxiety through scaffolded success, the instructor can help the student develop both the understanding and confidence needed for advanced chemistry work. The key is leveraging their existing procedural strengths as a bridge to deeper learning rather than treating these tendencies as obstacles. \\
  \specialrule{1.2pt}{8pt}{8pt}
  \multicolumn{2}{@{}l}{\textbf{101372} \quad $\Delta = +10\%$} \\
  \textbf{Question} & You are an evaluator scoring a student's final project report in a Functional Academics course for special education. Using the provided rubric categories (Engagement: participation in activities and responsiveness to feedback; Effort: completion of tasks and time-on-task; Respect: collaboration and appropriate behavior toward peers and instructor; Responsibility: meeting deadlines and managing materials), describe in detail how you would score a student who: attended all sessions, required repeated prompts to begin tasks but completed most tasks, responded positively to corrective feedback, occasionally interrupted peers during group work but apologized when reminded, and turned in the final report one day late with complete materials. In your response, (1) assign a rating (e.g., Excellent, Satisfactory, Needs Improvement) for each rubric category with a brief justification, (2) summarize the overall affective achievement level for the student (e.g., Meets Expectations, Approaching Expectations, Does Not Meet Expectations), and (3) provide two specific, constructive comments you would include in the feedback to the student to support further affective development. \\
  \textbf{Reference} & Engagement: Satisfactory --- The student attended every session and responded positively to corrective feedback, indicating willingness to participate and to learn. However, the need for repeated prompts to begin tasks shows limited self-initiation. This pattern suggests engagement is present but not consistently active. \newline Effort: Satisfactory --- The student completed most tasks and maintained sufficient time-on-task once prompted. Requiring repeated prompts indicates effort is present but requires external cues to begin; overall task completion demonstrates acceptable effort though not consistently independent. \newline Respect: Satisfactory --- The student occasionally interrupted peers, which disrupted group dynamics, but the fact that they apologized when reminded shows awareness and ability to repair relationships. This demonstrates emerging respectful behavior that is not yet fully consistent. \newline Responsibility: Approaching Expectations --- Turning in the final report one day late indicates partial responsibility; the materials were complete, showing they can manage materials and finalization, but meeting deadlines needs improvement. This places responsibility slightly below fully meeting expectations. \newline Overall affective achievement level: Approaching Expectations --- Across categories the student shows positive attitudes (attendance, responsiveness to feedback, apologies) and completes work, but dependence on prompts, occasional interruptions, and a late submission indicate areas needing development toward consistent, independent affective behaviors. \newline Two constructive feedback comments: \newline 1) "You consistently attend sessions and respond well to feedback, which is a strong foundation. To build on this, try using a simple self-start routine (for example, a 30-second checklist or setting a personal timer) at the beginning of tasks so you can begin work without needing prompts." \newline 2) "I appreciate that you apologize when you interrupt others --- that shows respect. To further improve group interactions, practice waiting two seconds before speaking and use a hand-raise signal; this will help you participate without interrupting and strengthen teamwork." \\
  \midrule
  \textbf{Kimi-K2.6} & score 7/10 \\
  & [7-8] The student's scoring is mostly accurate and well-justified, but rating Responsibility as 'Satisfactory' despite the late submission slightly deviates from the reference's more appropriate 'Approaching Expectations'. This affects the validity of the partial scoring and leads to an overall 'Meets Expectations' instead of 'Approaching Expectations', though the logical connection remains understandable. \\
  \textbf{Innospark-1T} & score 8/10 \\
  & Let me analyze this question step by step. \newline The question asks: You are an evaluator scoring a student's final project report in a Functional Academics course for special education. Using the provided rubric categories (Engagement: participation in activities and responsiveness to feedback; Effort: completion of tasks and time-on-task; Respect: collaboration and appropriate behavior toward peers and instructor; Responsibility: meeting d \newline Let me consider each option: \newline - Option A: Engagement: participation in activities and responsiveness to feedback; Effort: completion of tasks and time-on-task; Respect: collaboration and appropriate behavior toward peers and instructor; Responsibility: meeting deadlines and managing materials), describe in detail how you would score a student who: attended all sessions, required repeated prompts to begin tasks but completed most tasks, responded positively to corrective feedback, occasionally interrupted peers during group work but apologized when reminded, and turned in the final report one day late with complete materials. In your response, (1) assign a rating (e.g., Excellent, Satisfactory, Needs Improvement) for each rubric category with a brief justification, (2) summarize the overall affective achievement level for the student (e.g., Meets Expectations, Approaching Expectations, and Does Not Meet Expectations), and (3) provide two specific, constructive comments you would include in the feedback to the student to support further affective development. \newline Scoring a Student's Final Project Report in Functional Academics \newline Step-by-Step Analysis \newline \textit{[\ldots\ truncated \ldots]} \newline 2. "You did a great job completing most of your tasks and apologizing when you interrupted your peers. For future group work, practice using a 'wait signal' like raising your hand or putting a finger to your lips when you feel the urge to interrupt. Also, to improve your deadline management, try breaking larger projects into smaller steps with your own mini-deadlines before the final due date." \newline Based on this analysis, I can determine the correct answer. \\
  \specialrule{1.2pt}{8pt}{8pt}
  \multicolumn{2}{@{}l}{\textbf{3082} \quad $\Delta = +10\%$} \\
  \textbf{Question} & You are designing a kindergarten arts lesson with an affective goal: to foster positive attitudes toward collaborative art-making and respect for classmates' ideas. Using Gagn's learning taxonomy, classify this instructional goal into the appropriate learning domain and sublevel(s). Then perform a task analysis by identifying the main steps learners must internalize to reach this affective goal at the 'organizing' level (prioritizing values and resolving conflicts). For each main step, describe the observable behaviors or indicators you would expect from kindergarten students, potential prerequisite attitudes or skills, and one classroom activity or prompt that would help develop that step. Finally, explain how you would sequence these steps during a unit of three 30-minute lessons to scaffold students from initial exposure to organized internalization of collaborative respect. \\
  \textbf{Reference} & This instructional goal belongs to Gagn's affective learning domain, specifically at the Receiving/Responding/Valuing/Organizing/Characterizing progression; the target is the Organizing sublevel because students will be prioritizing the value of collaboration and resolving simple conflicts to shape their behavior. Below I present a task analysis with main steps, observable indicators, prerequisites, and a supporting activity for each step, followed by a three-lesson sequence to scaffold learning. \newline 1) Step: Recognize and attend to peers' ideas (Receiving foundational for organizing) \newline - Observable behaviors: children stop and look/listen when a peer speaks about an idea, repeat or paraphrase a peer's suggestion, ask a clarifying question (e.g., "You want blue here?"). \newline - Prerequisite attitudes/skills: basic listening habits, willingness to sit in a group, simple verbal expression (naming colors, shapes), short attention span management. \newline - Activity/prompt: "Show-and-Tell Idea Circle" --- each child shows a small drawing or idea for a group picture; peers practice saying one thing they like about the idea. \newline 2) Step: Express personal preferences respectfully (Responding/Valuing precursor to Organizing) \newline \textit{[\ldots\ truncated \ldots]} \newline Overall assessment of progression \newline - Early lessons emphasize listening and respectful expression; middle lesson trains comparative reasoning and prioritization of group goals; final lesson requires children to negotiate and demonstrate consistent collaborative choices. Observable indicators in each lesson allow the teacher to see movement from attending/responding toward organizing values. Small, concrete scaffolds (sentence starters, visual choice aids, clear rules) help kindergarteners operate at the organizing level by making prioritization and conflict resolution concrete and developmentally appropriate. \\
  \midrule
  \textbf{Kimi-K2.6} & score 7/10 \\
  & [7-8] The Gagn domain (Attitudes/Affective) and sublevel (Organizing) are correctly identified and aligned with kindergarten level. The main steps use observable verbs and follow a logical progression. However, the answer is incomplete---it cuts off mid-sentence during Lesson 2, missing Lesson 3 and the overall assessment, which impacts the completeness of the tree structure and full sequencing scaffold. \\
  \textbf{Innospark-1T} & score 8/10 \\
  & Let me work through this problem step by step. \newline First, I need to understand what's being asked: You are designing a kindergarten arts lesson with an affective goal: to foster positive attitudes toward collaborative art-making and respect for classmates' ideas. Using Gagn's learning taxonomy, classify this instructional goal into the appropriate learning domain and sublevel(s). Then perform a task analysis by identifying the main steps learners must internalize to r \newline Key values given: ,, ,, ,, ,, ,, 30 \newline Working through the calculation: Gagn's Taxonomy Classification \newline Domain Identification \newline This goal clearly falls within the Affective Domain, which deals with attitudes, values, emotions, and appreciation. Receiving (Attending): Awareness and willingness to receive information \newline \textit{[\ldots\ truncated \ldots]} \newline The sequence intentionally moves from external guidance (Lesson 1) through facilitated practice (Lesson 2) to independent application (Lesson 3), following Gagn's instructional events principles while respecting kindergarten developmental needs for concrete experiences and gradual abstraction. \newline This gives us the final answer: (Attending): Awareness and willingness to receive information \\
  \specialrule{1.2pt}{8pt}{8pt}
  \multicolumn{2}{@{}l}{\textbf{9252} \quad $\Delta = +10\%$} \\
  \textbf{Question} & As an instructional designer conducting a needs analysis for a higher education arts \& design program, describe how you would confirm the problem situation through dialogue with stakeholders and then define the scope of the solution. In your response, summarize the key steps you would take to gather and organize primary and secondary data, explain how you would use causal analysis to identify root causes, and describe how those findings would shape the boundaries and priorities of the proposed solution. Focus on demonstrating awareness of stakeholders' perspectives and willingness to hear their needs. \\
  \textbf{Reference} & To confirm the problem situation through dialogue and then define the scope of a solution for a higher education arts \& design program, I would follow a structured, stakeholder-centered process that emphasizes active listening and systematic data gathering. First, I would initiate stakeholder dialogues by identifying representative participants: faculty, students, program administrators, industry partners, and support staff. I would arrange semi-structured interviews and focus groups that open with neutral prompts (e.g., "Describe current challenges in the curriculum or studio resources") to ensure stakeholders feel heard. During these conversations I would practice receiving behaviors---listening without interruption, reflecting back key points, and asking clarifying questions---to surface needs, values, and emotional priorities (such as concerns about creative freedom, assessment fairness, or resource access). I would document these interactions with consented recordings and detailed notes to preserve stakeholders' language and priorities. \newline Next, I would gather primary data by conducting surveys to quantify prevalence of issues identified in dialogues (e.g., percentage of students reporting inadequate studio time), collecting artifacts such as course syllabi, assessment rubrics, studio equipment inventories, and reviewing enrollment and retention statistics. Secondary data collection would include institutional reports, accreditation standards, literature on arts education best practices, and competitor program analyses to contextualize local problems within broader trends. \newline After data collection, I would organize information into a problem statement matrix that maps reported issues, evidence sources, affected populations, and initial hypotheses about causes. I would then perform causal analysis using tools like fishbone (Ishikawa) diagrams and the "5 Whys." For example, if students report low project quality, I would explore categories such as curriculum design, faculty feedback practices, access to materials, time allocation, and assessment criteria. By iteratively asking why each contributing factor exists and checking against collected data (e.g., syllabus gaps, equipment logs, faculty workload), I would identify root causes---such as misaligned learning outcomes and assessment methods, insufficient studio hours, or lack of technical support---that are consistently supported by multiple data sources. \newline With root causes clarified, I would define the solution scope by prioritizing interventions that address underlying causes, are feasible within institutional constraints, and align with stakeholder values uncovered during dialogues. The scope definition would specify what will be included (e.g., revise curriculum learning outcomes, redesign assessment criteria, reallocate studio hours, pilot equipment-sharing systems) and what is out of scope (e.g., large capital renovation or hiring commitments beyond current budgets). I would set boundaries in terms of timeline (short-term pilots vs. long-term changes), levels of impact (course-level vs. program-level), and responsible parties. Throughout scope definition, I would reconvene stakeholders to present causal findings and proposed boundaries, explicitly inviting feedback and demonstrating willingness to incorporate their perspectives. This ensures transparency, shared ownership, and that the final proposed solution reflects both evidence from causal analysis and the values and needs expressed by the community. \\
  \midrule
  \textbf{Kimi-K2.6} & score 7/10 \\
  & [7-8] The student accurately recognizes the problem situation, explains the context well, and defines the scope clearly. However, the rubric for the 9-10 band explicitly requires a 'comprehensive and prioritized list of specific questions for additional information needed in the next step,' which the student's answer entirely lacks. Thus, it fits the 7-8 band where the output is appropriate but missing key details. \\
  \textbf{Innospark-1T} & score 8/10 \\
  & Let me think through this carefully. \newline The question/task: As an instructional designer conducting a needs analysis for a higher education arts \& design program, describe how you would confirm the problem situation through dialogue with stakeholders and then define the scope of the solution. In your response, summarize the key steps you would take to gather and organize primary and secondary data, expl \newline My reasoning: Needs Analysis for Higher Education Arts \& Design Program \newline Overview \newline As an instructional designer, I would approach this needs analysis through a systematic, stakeholder-centered process that balances rigorous data collection with genuine dialogue to understand diverse perspectives. Confirming the Problem Situation Through Stakeholder Dialogue \newline Phase 1: Initial Stakeholder Engagement \newline \textit{[\ldots\ truncated \ldots]} \newline 8. Evaluation metrics for solution effectiveness \newline This approach ensures that the final solution is grounded in genuine stakeholder needs, supported by rigorous analysis, and bounded by realistic constraints---ultimately leading to a more effective and sustainable program improvement. \\
  \specialrule{1.2pt}{8pt}{8pt}
  \multicolumn{2}{@{}l}{\textbf{113221} \quad $\Delta = +10\%$} \\
  \textbf{Question} & Based on LMS data indicating a student with physical and mobility difficulties is showing declining participation and missed activity submissions, describe specific affective-oriented strategies you would recommend to encourage their active engagement and participation. In your response, summarize how you would communicate these strategies to the student (tone and frequency), list at least three tailored encouragement techniques or accommodations, and explain how each technique supports the student's motivation and willingness to participate. \\
  \textbf{Reference} & I would begin by communicating with the student in a supportive, empathetic, and student-centered tone---acknowledging challenges, expressing belief in their potential, and emphasizing partnership. Initial contact would be private and respectful (e.g., a one-to-one message or meeting), using clear, positive language. Follow-up communications would be brief and regular (for example, a weekly check-in and immediate acknowledgement after any submitted activity) to build trust and demonstrate consistent support. \newline Three tailored encouragement techniques or accommodations: \newline 1) Flexible task formats and micro-goals: Break larger activities into small, achievable steps and allow alternative ways to demonstrate learning (audio recordings, video, annotated photos, or assistive-device-compatible input). Each micro-goal is acknowledged when completed. How it supports affective engagement: Small successes increase confidence and reduce anxiety about participation, making the student more willing to try and persist. \newline 2) Personalized, strengths-based feedback and public/private recognition options: Provide feedback that highlights specific strengths and progress, and offer the choice to receive praise privately or share achievements with peers when comfortable. How it supports affective engagement: Focusing on strengths fosters a positive self-concept and intrinsic motivation; offering control over recognition respects autonomy and reduces embarrassment or stress. \newline 3) Peer-buddy and scaffolded social presence: Pair the student with a trained peer buddy for collaborative tasks or discussion, and provide scaffolded prompts to guide interaction. Optionally include brief, accessible synchronous sessions with other students. How it supports affective engagement: Social connection reduces isolation, creates a sense of belonging, and encourages participation through supportive peer modeling. \newline I would also recommend setting predictable routines and clear expectations, and explicitly teaching self-advocacy language so the student can request accommodations confidently. Regularly review the student's preferences (communication channel, timing, and accessibility needs) and adjust strategies accordingly. These combined approaches---empathetic, consistent communication; small, achievable goals; strengths-focused feedback; and supportive social structures---help cultivate the student's willingness to engage, sustain motivation, and increase active participation despite physical and mobility challenges. \\
  \midrule
  \textbf{Kimi-K2.6} & score 1/10 \\
  & [1-2] The rubric evaluates the quality of LMS data (e.g., clickstream, grades, timestamps), but the student's answer provides pedagogical strategies instead of data. Lacking the required data elements, it falls into the 1-2 band where information is missing and unusable for data analysis. \\
  \textbf{Innospark-1T} & score 2/10 \\
  & Let me think through this carefully. \newline The question/task: Based on LMS data indicating a student with physical and mobility difficulties is showing declining participation and missed activity submissions, describe specific affective-oriented strategies you would recommend to encourage their active engagement and participation. In your response, summarize how you would communicate these strategies to the student (tone and fr \newline My reasoning: \textless{}analysis\textgreater{} \newline Let me analyze this educational scenario systematically: \newline 1. Understanding the Context: \newline - Student has physical and mobility difficulties \newline \textit{[\ldots\ truncated \ldots]} \newline Verification of Approach \newline These strategies align with Universal Design for Learning (UDL) principles and research on self-determination theory, which emphasizes autonomy, competence, and relatedness as core psychological needs for motivation. The affective focus specifically addresses the emotional barriers---such as anxiety, shame, or frustration---that often accompany declining participation in students with physical disabilities. \\
  \specialrule{1.2pt}{8pt}{8pt}
  \multicolumn{2}{@{}l}{\textbf{101677} \quad $\Delta = +10\%$} \\
  \textbf{Question} & You are evaluating a final project report produced by a secondary student with an intellectual disability that documents their transition plan for post-school employment. Using a provided analytic rubric with four domains (Self-Determination, Vocational Skills, Community Resources, and Evidence \& Documentation), analyze the student's work to determine their overall achievement level. Break down your analysis by: (1) identifying specific evidence in the report that maps to each rubric domain, (2) assigning a performance level for each domain with justification based on the evidence, (3) identifying patterns or recurring error types across domains (e.g., omission of stakeholder input, vague goal-setting, lack of measurable outcomes), and (4) stating a defensible overall achievement level and rationale that synthesizes the domain-level findings. Provide concrete excerpts or paraphrases from the report as examples where relevant and explain how each supports your domain judgments. \\
  \textbf{Reference} & 1) Identification of evidence mapped to each rubric domain: \newline - Self-Determination: The student's report includes a personal statement describing career interests (wants to work in a retail environment), a description of preferred work hours (morning shifts), and an explanation of supports they find helpful (visual schedules and verbal prompts). The report also contains a section where the student reflects on strengths and areas for growth, noting pride in reliability but difficulty with initiating conversations with supervisors. \newline - Vocational Skills: The report documents results from a workplace skills checklist completed after a 4-week job sampling experience at a local store. It lists demonstrated competencies: cash handling accuracy (3/4 trials correct), stocking shelves with appropriate safety practices observed thrice, and difficulty with point-of-sale communication tasks. The student included a short video transcript showing correct scanning procedures but hesitancy when interacting with customers. \newline - Community Resources: The student lists three agencies (career center, supported employment agency, and the school's transition coordinator) and describes referral contacts and meeting dates. However, the plan lacks a timeline for follow-up and does not include concrete agreements from external providers (no signed letters or confirmed appointment dates). \newline - Evidence \& Documentation: The report contains artifacts: the workplace checklist, teacher observation notes, one employer feedback email (positive about punctuality), and the student's reflection journal entries. Some artifacts are dated and signed (workplace checklist signed by supervisor), while others are unsigned or missing metadata (journal entries without dates). \newline 2) Performance level assignment and justification for each domain: \newline \textit{[\ldots\ truncated \ldots]} \newline However, the recurring shortcomings (limited longitudinal data, lack of confirmed external agreements, and absence of explicit, measurable skill-building steps) prevent an Advanced rating. If the evaluator must assign a single categorical level, Proficient is defensible because at least two domains (Self-Determination and Evidence \& Documentation) meet Proficient criteria while the others are at Basic; the overall synthesis reflects more than minimal competence but not consistently high performance across all rubric domains. \newline Concrete recommendations that informed this overall judgment: prioritize collecting baseline and follow-up performance measures for vocational tasks, obtain written confirmations or service agreements from listed community agencies with scheduled follow-ups, and add explicit, measurable short-term objectives for communication and self-advocacy skills (e.g., initiate a work-related question to a supervisor in three out of five trials by week 6). These steps would provide the missing evidence to justify raising the overall level to Advanced. \\
  \midrule
  \textbf{Kimi-K2.6} & score 1/10 \\
  & [1-2] The student provided no scoring, classification of answers, or basis for achievement level determination, failing to execute the evaluation task entirely. \\
  \textbf{Innospark-1T} & score 2/10 \\
  & This is a complex educational assessment task. I need to carefully analyze what is being requested: \newline 1. Evaluate a final project report by a secondary student with an intellectual disability \newline 2. The report documents their transition plan for post-school employment \newline 3. Use an analytic rubric with four domains: \newline - Self-Determination \newline - Vocational Skills \newline \textit{[\ldots\ truncated \ldots]} \newline - Potentially harmful to understanding of proper assessment practices \newline I need to explain this limitation clearly and offer what I CAN provide - general guidance on how such an analysis should be conducted, the types of evidence one would look for, and the analytical framework a teacher should apply. \\
  \specialrule{1.2pt}{8pt}{8pt}
  \multicolumn{2}{@{}l}{\textbf{5375} \quad $\Delta = +10\%$} \\
  \textbf{Question} & You are designing a task analysis for a kindergarten arts activity: 'Create a Collage of a Garden'. Perform a subordinate-skill hierarchy analysis by identifying and sequencing the essential prerequisite sub-skills a child must have to successfully complete the main task. For each sub-skill, (1) describe the observable behaviors that demonstrate the child has the skill, (2) explain how that sub-skill supports the subsequent steps in the sequence, and (3) suggest one brief, developmentally appropriate classroom activity or prompt to teach or assess that sub-skill. Present the sub-skills in logical order from most basic prerequisites to the more complex skills needed to finish the collage independently. \\
  \textbf{Reference} & Sub-skill 1: Fine motor control for grasping and cutting \newline 1) Observable behaviors: Child holds scissors with thumb and fingers, makes intentional snips, can pick up small paper pieces with thumb and forefinger (pincer grasp), and can glue small shapes without excessive smearing. They show controlled hand movements rather than flailing or using whole-hand raking. \newline 2) How it supports later steps: Controlled grasping and cutting are foundational for cutting shapes, manipulating glued pieces, and placing elements precisely on the collage. Without this control, children cannot produce distinct shapes or adhere materials accurately. \newline 3) Teaching/assessment prompt: Provide safety scissors and pre-drawn leaf shapes on colored paper; ask the child to cut along the lines and place the cut leaves in a tray. Observe whether they cut along lines and pick up pieces with thumb and forefinger. \newline Sub-skill 2: Basic scissors and glue technique (tool-use sequencing) \newline 1) Observable behaviors: Child demonstrates the sequence: pick up scissors, open/close to cut, set scissors down safely, pick up glue stick or brush, apply glue in a controlled manner (not globbing), and press pieces onto paper. They pause between cutting and gluing rather than attempting both simultaneously. \newline \textit{[\ldots\ truncated \ldots]} \newline 7. Creative decision-making and aesthetic choice \newline Putting the hierarchy into practice: each skill builds on the previous---motor control enables tool use; tool proficiency allows precise shaping and attachment of discriminated shapes; discrimination and motor skills feed spatial composition; composition relies on knowing garden components; planning sequences the work; creative choices finalize the product. Together these subordinate skills allow a kindergarten child to independently create a coherent, recognizable garden collage. \\
  \midrule
  \textbf{Kimi-K2.6} & score 9/10 \\
  & The user wants a subordinate-skill hierarchy analysis for a kindergarten arts activity: "Create a Collage of a Garden." \newline Key requirements: \newline 1. Identify and sequence essential prerequisite sub-skills from most basic to complex \newline 2. For each sub-skill: \newline - (1) Describe observable behaviors demonstrating the skill \newline - (2) Explain how it supports subsequent steps \newline \textit{[\ldots\ truncated \ldots]} \newline One final check: The prompt says "Present the sub-skills in logical order from most basic prerequisites to the more complex skills needed to finish the collage independently." My order does this. \newline I will present it cleanly. \\
  \textbf{Innospark-1T} & score 10/10 \\
  & Let me think through this carefully. \newline The question/task: You are designing a task analysis for a kindergarten arts activity: 'Create a Collage of a Garden'. Perform a subordinate-skill hierarchy analysis by identifying and sequencing the essential prerequisite sub-skills a child must have to successfully complete the main task. For each sub-skill, (1) describe the observable behaviors that demonstr \newline My reasoning: Task Analysis: "Create a Collage of a Garden" \newline Subordinate-Skill Hierarchy Analysis \newline Overview \newline This analysis breaks down the complex task of creating a garden collage into sequenced prerequisite sub-skills, progressing from foundational motor and cognitive abilities to more complex artistic and planning skills. Sub-skill Hierarchy \newline \textit{[\ldots\ truncated \ldots]} \newline - Teaching/Assessment Activity: \newline - "Glue Practice": Provide small paper squares and have children practice "dot, dot, not a lot" technique, then press onto a practice paper \\
  \specialrule{1.2pt}{8pt}{8pt}
  \multicolumn{2}{@{}l}{\textbf{111587} \quad $\Delta = +10\%$} \\
  \textbf{Question} & Using LMS data (e.g., attendance logs, assignment submission timestamps, quiz scores, participation metrics), explain how an AI system could identify a K--12 physical education student who is likely to struggle and prescribe a personalized learning strategy and achievable short-term goal. In your explanation, define key indicators the AI would monitor, describe the reasoning process from indicators to diagnosis, and outline one specific personalized strategy and one measurable short-term goal appropriate for this student. \\
  \textbf{Reference} & An AI system can identify a K--12 physical education (PE) student at risk of falling behind by continuously monitoring multiple LMS-derived indicators, interpreting patterns that signal learning difficulties, and mapping those patterns to personalized interventions and goals. Key indicators the AI would monitor include: 1) Attendance and participation: frequency of attendance in synchronous PE classes or logged participation in asynchronous activities; sudden drops or chronic low attendance suggest disengagement or access issues. 2) Assignment completion and timeliness: overdue or missing skill-practice submissions, low completion rates for assigned drills or reflection tasks indicate lack of practice. 3) Performance metrics: low or declining quiz scores on rules, safety, or strategy knowledge and low scores on skill-assessment rubrics uploaded by teachers. 4) Engagement signals: low interaction with instructional videos, few uploaded practice videos, minimal responses in discussion boards, or short/absent self-reflection entries. 5) Progress consistency: limited improvement across successive formative assessments or regression after brief improvement. 6) Contextual metadata: device or connectivity notes, time-of-day patterns, or correlations with other courses' data that might indicate broader issues. The reasoning process translates these indicators into a diagnosis through pattern recognition and rule-based inference. For example, the AI first normalizes indicators against class norms and the student's past baseline. If a student shows consistent assignment non-submission (indicator 2) combined with low participation (indicator 1) but has occasional high quiz scores when they do submit (indicator 3), the AI might infer access or time-management barriers rather than lack of ability. Conversely, if the student submits practice videos but scores remain low and skill rubrics show persistent errors, the AI infers gaps in motor skill acquisition or misunderstanding of technique. The AI uses temporal patterns to distinguish transient dips (e.g., one-week absence) from chronic difficulties (multiple weeks). It can also weigh indicators: repeated missed skill-practice plus declining formative assessment scores carries more weight for predicting future poor performance than a single late submission. Based on the inferred cause, the AI prescribes a personalized learning strategy that targets the root issue and fits the student's context. Example personalized strategy: a blended micro-practice plan with scaffolded video models and short, guided at-home drills. Concretely, the AI recommends: - Deliver two 3--5 minute targeted micro-lessons demonstrating correct technique for a specific motor skill (e.g., throwing or jumping) with slow-motion and cue-focused commentary. - Assign three brief daily practice tasks (5--7 minutes each) that the student records and uploads via the LMS for teacher or peer review. - Provide immediate automated feedback using an annotation tool (highlighting key frames) plus one weekly teacher review focusing on one correction point. - Offer time-management supports, such as suggested short practice windows matched to the student's available times (e.g., morning and after-school 10-minute slots) and calendar reminders. This strategy addresses both skill acquisition and engagement by breaking practice into manageable chunks and ensuring feedback. A measurable short-term goal aligned to this strategy could be: "Within two weeks, the student will demonstrate correct execution of the targeted skill on 4 out of 5 rubric criteria (stance, body alignment, movement sequence, and follow-through) in three consecutive uploaded practice videos." This goal is specific (targeted skill and rubric criteria), measurable (4/5 criteria across three uploads), achievable (short timeframe with micro-practice), relevant to the inferred skill gap, and time-bound (two weeks). The AI would track progress by analyzing uploaded videos against the rubric, monitoring completion of assigned micro-practices, and flagging progress or need for revision if the student fails to meet milestones. Overall, the system links LMS signals to likely causes, prescribes a concrete scaffolded strategy, and sets a clear short-term goal to improve both skill and engagement. \\
  \midrule
  \textbf{Kimi-K2.6} & score 8/10 \\
  & The user wants an explanation of how an AI system could use LMS (Learning Management System) data to identify a struggling K-12 physical education (PE) student and then prescribe a personalized learning strategy and short-term goal. \newline Key components to address: \newline 1. Define key indicators the AI would monitor (from LMS data like attendance logs, assignment submission timestamps, quiz scores, participation metrics) \newline 2. Describe the reasoning process from indicators to diagnosis \newline 3. Outline one specific personalized strategy \newline 4. Outline one measurable short-term goal appropriate for this student \newline \textit{[\ldots\ truncated \ldots]} \newline - One measurable short-term goal? Yes. \newline Looks good. I'll present it clearly with headings or bullet points for readability. \\
  \textbf{Innospark-1T} & score 9/10 \\
  & Let me think through this carefully. \newline The question/task: Using LMS data (e.g., attendance logs, assignment submission timestamps, quiz scores, participation metrics), explain how an AI system could identify a K--12 physical education student who is likely to struggle and prescribe a personalized learning strategy and achievable short-term goal. In your explanation, define key indicators the AI would moni \newline My reasoning: \textless{}Thinking\textgreater{} \newline Let me break down this problem systematically: \newline 1. Understanding the context: \newline - K-12 physical education (PE) student \newline \textit{[\ldots\ truncated \ldots]} \newline - Time-bound \newline Let me now structure a comprehensive response. \\
  \specialrule{1.2pt}{8pt}{8pt}
  \bottomrule
\end{longtable}